\section{关系证据与能力证据：申菱环境、科大讯飞、拓维信息、航天电器}

\subsection{申菱环境（301018）：最接近“客户+产品+经营增速”，但仍差平台订单}

申菱环境是本案上游映射中公开证据最完整的公司。2025 年交易所完整年报原件列出华为为数据服务客户，并称与 H 公司、字节、腾讯、阿里及主要 COLO 客户合作粘性增强；数据服务板块收入同比增长 51.42\%，新增订单同比增长约 72\%。产品覆盖 CDU、Manifold、一体化冷源、干冷器、预制化管网、液冷门、快速接头和冷板，且新基建领域智能温控基地已经投产。以上信息同时给出客户、产品、收入增速、订单和产能方向，显著强于普通概念映射。

边界也同样清楚：年报没有出现 Atlas 950、UnifiedBus 或 950DT，没有平台认证、合同金额、CDU/冷板/快接数量、单价、毛利率和验收节奏。“H 公司”与公开客户华为可以支持一般合作，但不能推定具体项目。公司 FY2025 营收 42.09 亿元、归母净利 2.17 亿元，分别增长 39.55\%和 87.59\%；2026Q1 收入和归母净利同比下降 1.80\%和 47.71\%，说明订单高增尚未稳定转化为季度利润。券商 2026E EPS 表内均值为 1.608 元，但与现股本口径不勾稽；按 2026E 归母净利 4.28 亿元和现股本重算为 1.144 元，对应约 76 倍 PE。成熟业务 PE 与基于 FY2025 官方净资产的合理 P/B 双方法综合目标 31.50 元，低于现价 63.8\%。

投资含义：关系证据优先级最高，但价格已经给出较强兑现要求。升级触发是公司正式点名 Atlas 平台、液冷组件范围、订单/交付及盈利；下调触发是数据服务订单不能转为收入、应收/存货扩张或毛利率继续承压。

\subsection{科大讯飞（002230）：最强 950 下游协同证据，不是硬件供应链标的}

科大讯飞 2025 年年报明确披露与华为围绕昇腾 950 协同开发，并计划在 2026 年 10 月发布旗舰模型。这是本案最强的 950 系列下游工作负载证据，说明国产大模型厂商正在适配下一代算力底座。其价值在于验证需求与软件生态，而非证明公司取得上游订单或采购了多少 Atlas 950。

财务上，FY2025 营收 271.05 亿元、归母净利 8.39 亿元，但扣非净利仅 2.64 亿元；2026Q1 归母亏损 1.70 亿元、扣非亏损 4.30 亿元。2026H1 预告仍为 1.80--2.28 亿元亏损，只是亏损收窄。按 2026E 归母净利和现股本重算 EPS 约 0.503 元，对应约 81.8 倍 PE。即使模型发布与商业化如期推进，高估值仍要求付费用户、收入增速、销售效率和现金流同步改善。综合目标 16.01 元，低于现价 61.1\%，因此定位为“下游事件验证”，不作为 Atlas 硬件受益股。

\subsection{拓维信息（002261）：关系直接，盈利质量和外部预测门槛未过}

拓维信息不是没有产业链内容，而是需要把“生态关系”和“可估值收入”拆开。公司年报披露其为首批昇腾整机伙伴，子公司湘江鲲鹏取得昇腾部件伙伴认证，并以“兆瀚”品牌推出 RH220T-C 高性能服务器、RA2300-C 推理服务器、DeepSeek 一体机及边缘设备；公司参与长沙、重庆、济南等人工智能创新中心和多个政企算力项目。FY2025 对华为销售占比 41.51\%、采购占比 50.15\%，这确实是较强的华为/昇腾生态关系证据。

但上述证据仍停留在生态与整机/集成层。年报没有出现 Atlas 950、950DT、UnifiedBus、具体平台认证编号、客户定点、机柜数量、ASP、交付节奏或分部收入拆分；41.51\% 是全公司对华为销售集中度，不是 Atlas 950 收入占比。换言之，拓维可以证明“能做昇腾相关整机和项目”，不能证明“已拿到 Atlas 950 超节点订单”。

盈利质量解释了为什么本报告此前没有给它目标价：FY2025 营收 31.71 亿元、同比下降 22.79\%，归母净利 0.64 亿元，但扣非净利为 -0.40 亿元；经营活动现金流净额 7.78 亿元，现金回收尚可，却不能替代经常性利润。2026Q1 收入 5.94 亿元、同比下降 4.72\%，归母净利 0.61 亿元、同比下降 5.96\%，扣非净利仍为 -0.08 亿元。公司披露母公司未弥补亏损，分红条件亦未满足，说明利润改善尚未形成稳定股东回报能力。

估值上，180 天窗口内没有找到新鲜原始卖方覆盖；民生证券 2024 年历史报告仅给出 2026E 收入 57.20 亿元、归母净利 1.84 亿元和 EPS 0.15 元，没有 2027E 和明示目标价，不能作为当前目标价。我们的处置是“关系强、估值零权重”：在出现新鲜可审计的 2026E/2027E 预测、扣非持续转正、具体 Atlas 产品/订单与收入拆分前，不发布目标价。

升级触发应同时满足四项：华为/客户或公司点名 Atlas 950/950DT 产品与认证范围；披露订单、交付或客户验收；披露算力产品收入、毛利和应收/存货变化；扣非净利连续两个报告期转正。任何单项问答、合作伙伴称号或项目中标都不能单独把拓维从观察池升级。

\subsection{航天电器（002025）：高速背板与液冷连接器能力，客户关系和盈利经济性仍未闭环}

航天电器应当单独分析，而不能被一句“连接器能力观察”带过。公司核心产品覆盖高速背板连接器、高速连接器组件、液冷连接器和液冷管路，并在 2025 年“2+N”战略中把大互连、大电机与新域新质、战略性新兴产业、AI 算力、商业航天等并列推进。2026 年 4 月 20 日投资者问答中，公司确认“有算力领域用高速背板连接器、高速连接器组件、液冷连接器和液冷管路等产品，并为下游客户配套”，但没有确认提问前提中的华为昇腾 950 供货、客户名称、订单金额或营业额。因此，产品能力可以进入能力卫星池，华为/Atlas 专属收入仍必须记为 0。

基本面并不差，但利润和现金流尚未匹配主题估值。FY2025 营收 58.20 亿元、同比增长 15.82\%，归母净利 1.83 亿元、同比下降 47.32\%，营业利润下降 42.42\%；连接器及互连一体化产品收入 39.84 亿元，占总收入 68.46\%，同比增长 17.59\%，毛利率 29.07\%，同比下降 9.61 个百分点。经营活动现金流净额为 -6.14 亿元，较上年 -2.59 亿元进一步恶化，公司解释为订单增长下的采购与外协结算压力。2026Q1 收入 16.14 亿元、同比增长 10.09\%，归母净利 0.52 亿元、同比增长 12.18\%，扣非净利 0.45 亿元；改善仍需连续季度验证。

公司客户集中度和业务属性也要求折价：前五大客户合计占 2025 年销售额 49.17\%，第一大客户为航天科工下属关联企业，占比 21.07\%；研发投入 7.38 亿元，同比下降 1.75\%。公司给出的 2026 年预算是收入 65.00 亿元、成本及费用不超过 61.12 亿元，属于经营目标而非订单或盈利承诺。若高端互连放量而毛利率、应收、存货和现金流不能同步改善，收入增长并不能支撑高倍数。

外部预测分歧是本股的核心估值风险。4 份原始卖方报告的 2026E EPS 区间为 0.53--1.10 元、2027E 为 0.62--2.15 元；最新国信证券报告给出 69--78 元目标区间及 2026E/2027E EPS 0.85/1.18 元。按本报告统一现股本重算，2026E EPS 为 0.836 元，现价 64.76 元对应约 77.4 倍 PE；成熟业务 PE、合理 P/B 与外部目标锚综合目标为 42.06 元，较现价低 35.1\%。该目标评价的是公司更广义的连接器/军工业务，并不包含 Atlas 950 增量。

投资动作是“能力观察、估值不追”：只有当公司或客户正式披露 950 平台资格、具体背板/液冷 SKU、订单与交付、收入和毛利桥后，才上调 Atlas 相关信用；若 2026 年收入预算落空、连接器毛利率继续下降、经营现金流转负或大客户集中度上升，则先下调广义估值。航天电器是可验证的高速互连能力标的，不是已确认的华为 Atlas 950 供应商。

\section{广义 AIDC 盈利基础：光互联、PCB/CCL 与热管理}

\subsection{华工科技（000988）：强光连接盈利基础，Atlas 关系为零信用期权}

华工科技 2025 年光电器件收入 60.97 亿元，同比增长 53.39\%，覆盖 800G、1.6T、3.2T 光引擎及 AEC/ACC 等产品；FY2025 总营收 143.55 亿元、归母净利 14.71 亿元，2026Q1 归母净利同比增长 55.76\%。这构成可审计的广义 AI 光互联盈利基础。问题在于 Atlas 950 的端口、模块和供应商均未披露，且 UBoE/\allowbreak CPO/\allowbreak LPO/\allowbreak OCS 路线可能改变模块价值量。

按现股本重算的 2026E EPS 为 2.361 元，对应约 49.8 倍 PE。综合目标 55.48 元，低于现价 52.8\%。公司应被视为“AI 光连接强基本面+Atlas 可选性”，而非已确认华为新平台供应商。

\subsection{英维克（002837）与科华数据（002335）：能力完整，项目边界决定利润}

英维克拥有从冷板、快接、Manifold、CDU 到冷源、工质和测试设备的端到端液冷能力，年报披露部分主流算力芯片厂商和头部设备厂商认可并开始规模采购；FY2025 营收 60.68 亿元、归母净利 5.22 亿元。但 2026Q1 归母净利同比下降 81.97\%，表明收入增长不能替代毛利和季度交付验证。券商表内 2026E EPS 1.171 元存在送转股口径差异；按现股本重算为 0.897 元、约 68.6 倍 PE，综合目标 23.64 元，低于现价 61.6\%。

科华数据拥有智能算力中心供配电、200kW--1.5MW 液冷机组和 40--120kW/柜 POD 能力，并披露北京昇腾创新中心相关案例。FY2025 归母净利增长 32.62\%，2026H1 预告归母净利增长 50\%--60\%，但扣非增速仅约 2.79\%--13.68\%，需要辨别非经常性因素和项目验收质量。券商表内 2026E EPS 1.35 元与现股本口径不勾稽；重算 EPS 为 0.921 元、约 32.3 倍 PE，综合目标 15.18 元，低于现价 49.0\%。其 Atlas 专属性也弱，定位为“等待估值回落”。

\subsection{胜宏科技、沪电股份、深南电路、生益科技：景气真实，专属关系缺失}

AI PCB/CCL 是 2025--2026 年财务兑现最强的广义 AIDC 环节之一。胜宏科技 FY2025 营收增长 79.77\%、归母净利增长 273.52\%；沪电股份 2026H1 预告归母净利增长约 68\%--78\%；深南电路预告增长约 54\%--69\%；生益科技预告增长约 117\%--131\%。这些数字证明高端板材、AI 服务器/交换机 PCB 和产品结构升级正在形成利润。

然而，Atlas 950 的 PCB 层数、材料体系、供应商认证和分配均未披露。当前保守模型中，胜宏科技、沪电股份、深南电路和生益科技的综合目标空间分别约为 -37.0\%、-51.5\%、-39.8\%和 -48.0\%。该组适合按广义 AI 订单、良率、产品结构和估值选择，不应以华为发布会作为唯一催化。

\section{低纯度价值卫星：沃尔核材与连接器公司}

沃尔核材按现股本重算的 2026E EPS 为 1.343 元，现价对应约 11.3 倍 PE，综合目标 15.45 元、空间仅 +1.4\%。其吸引力来自较低估值和高速铜连接相关业务预期，不来自已确认的 Atlas 950 订单。若后续无法获得平台认证、收入拆分和产能利用率证据，估值修复只能依赖公司原有业务和广义高速连接景气。

兆龙互连、鼎通科技等同样进入能力池，但没有完成 Atlas 供应链证据闭环；航天电器的完整公司卡见上文。

\clearpage
\begin{exhibitbox}[重点公司财务与估值快照]
\scriptsize
\setlength{\tabcolsep}{1pt}
\begin{tabularx}{\linewidth}{L{1.5cm}L{2.1cm}R{1.7cm}R{1.7cm}R{1.7cm}R{1.7cm}R{1.7cm}X}
\toprule
\textbf{代码} & \textbf{公司} & \textbf{2025 收入} & \textbf{2025 归母} & \textbf{2026E EPS} & \textbf{2026E PE} & \textbf{目标空间} & \textbf{Atlas 信用}\\
\midrule
301018 & 申菱环境 & 42.09 & 2.17 & 1.144 & 76.0x & -63.8\% & 关系最强；专属收入 0\\
000988 & 华工科技 & 143.55 & 14.71 & 2.361 & 49.8x & -52.8\% & 光连接能力；专属收入 0\\
002837 & 英维克 & 60.68 & 5.22 & 0.897 & 68.6x & -61.6\% & 液冷能力；专属收入 0\\
002335 & 科华数据 & 81.60 & 4.18 & 0.921 & 32.3x & -49.0\% & 昇腾邻近项目；专属收入 0\\
002130 & 沃尔核材 & 84.51 & 11.44 & 1.343 & 11.3x & +1.4\% & 高速连接能力；关系未确认\\
002230 & 科大讯飞 & 271.05 & 8.39 & 0.503 & 81.8x & -61.1\% & 950 下游协同；硬件收入 0\\
002025 & 航天电器 & 58.20 & 1.83 & 0.836 & 77.4x & -35.1\% & 背板/液冷连接器；950 未确认\\
002261 & 拓维信息 & 31.71 & 0.64 & -- & -- & -- & 昇腾整机生态；不估值\\
300476 & 胜宏科技 & 192.92 & 43.12 & 9.054 & 26.7x & -37.0\% & AI PCB；Atlas 关系未确认\\
600183 & 生益科技 & 284.31 & 33.34 & 2.292 & 57.7x & -48.0\% & AI CCL；Atlas 规格未知\\
\bottomrule
\end{tabularx}
\sourcenote{单位：亿元，EPS 为元。12 家公司交易所年报原件、原始卖方盈利预测与现股本重算模型。}
\end{exhibitbox}

\section{公司卡的共同审计问题}

未来每家公司都用同一组六问复核：第一，是否点名 Atlas/UnifiedBus/950 平台与具体产品；第二，是否完成认证、定点或订单；第三，是否披露在手订单、交付和收入；第四，毛利率、ASP 与客户集中度是否支持利润；第五，存货、应收和资本开支是否把会计利润转成现金；第六，当前价格是否已经提前计入多年增长。缺少前三问时，Atlas EPS 必须为 0；缺少后三问时，即使关系确认也不能给高倍数。
